\documentclass[ 12pt,a4paper,twocolumn,fleqn]{report}
\usepackage{graphicx}
\usepackage[a4paper,top=20mm, bottom=30mm, left=10mm, right=25mm]{geometry}
\usepackage{fancyhdr}
\usepackage{hyperref}
\usepackage{listings}
\usepackage[]{algorithm2e}
\usepackage{color}
\usepackage{fancybox}
\usepackage{lastpage}
\thisfancypage{%
  \setlength{\fboxsep}{20pt}\doublebox}{}
\pagestyle{fancy}
\usepackage{lineno}
\usepackage{xtab,booktabs}
\usepackage{setspace}
\usepackage{amsmath}
\usepackage{xcolor}
\usepackage{tocloft} 
\addtolength{\cftaftertoctitleskip}{-\baselineskip} 
\usepackage[compact]{titlesec}
\titlespacing{\section}{0pt}{*0}{*0}
\titlespacing{\subsection}{0pt}{*0}{*0}
\titlespacing{\subsubsection}{0pt}{*0}{*0}
\setlength\columnsep{25pt}
\makeatletter
\g@addto@macro{\normalsize}{%
\setlength{\abovedisplayskip}{0pt}
\setlength{\abovedisplayshortskip}{0pt}
\setlength{\belowdisplayskip}{0pt}
\setlength{\belowdisplayshortskip}{0pt}}
\makeatother
\mathindent=0.0pt
\usepackage{float}
\renewcommand{\baselinestretch}{1.0}
\pagenumbering{global}
\fancypagestyle{myfooter}{
    \fancyhf{} 
    \fancyfoot[L]{Department of CSE} 
    \fancyfoot[R]{\protect\thepage} 
    \renewcommand{\headrulewidth}{0pt} 
    \renewcommand{\footrulewidth}{0.4pt} 
}

% Redefine plain page style for chapters
\fancypagestyle{plain}{
    \fancyhf{} 
    \fancyfoot[R]{\protect\thepage} 
    \renewcommand{\headrulewidth}{0pt} 
    \renewcommand{\footrulewidth}{0pt} 
}

% Define a new page style combining fancy and myfooter
\fancypagestyle{combined}{
    \fancyhf{} 
    \fancyhead{} 
    \fancyhead[L]{Crop Stress Grading using Hyperspectral Data and Deep Learning} 
    \fancyfoot[L]{Department of CSE (AI \& ML)} 
    \fancyfoot[R]{\protect\thepage} 
    \renewcommand{\headrulewidth}{0.4pt} 
    \renewcommand{\footrulewidth}{0pt} 
}
\begin{document}

\fancyfoot[R]{Page \thepage}  
\pagestyle{plain}

\onehalfspacing
\onecolumn
\begin{center}
    \huge\textbf{DAYANANDA SAGAR UNIVERSITY} \\
    \large\textbf{Devarakaggalahalli, Harohalli,
Kanakapura Road,
Bengaluru South Dt. – 562 112 }
\end{center}
\begin{figure}[h]
    \centering
    \includegraphics[width=0.4\textwidth]{image/SCHOOL OF ENGINEERING_c.png}
\end{figure}
\begin{center}
    \Large\textbf{Bachelor of Technology} \\
    \normalsize\textbf{in} \\
    \large\textbf{COMPUTER SCIENCE AND ENGINEERING} \\
    \large\textbf{(ARTIFICIAL INTELLIGENCE AND MACHINE LEARNING)}
\end{center}
\vspace{10mm}
\begin{center}
    \large\textbf{A Project Report On} \\
    \large\textbf{ Crop Stress Grading using Hyperspectral Data and Deep Learning}
\end{center}
\vspace{10mm}
\begin{center}
    \small\textbf{By} \\
    \normalsize\textbf{VANSH SONI} \\
\end{center}
\begin{figure}[h]
    \centering
    \includegraphics[width=0.11\textwidth]{image/Logo of AIML.jpg}
\end{figure}
\begin{center}
    \normalsize\textbf{Under the supervision of} \\
    \normalsize\textbf{Prof. UDAYA BHASKARA N} \\
    \normalsize\textbf{Assistant Professor}\\
    \normalsize\textbf{Computer Science \& Engineering (AI \& ML)}
\end{center}
\vfill
\begin{center}
    \normalsize\textbf{SCHOOL OF ENGINEERING} \\
\end{center}



\newpage
  \pagestyle{plain}
\thisfancypage{%
  \setlength{\fboxsep}{20pt}\doublebox}{}
\begin{center}
\textcolor{blue}{\LARGE\textbf{{DAYANANDA SAGAR UNIVERSITY}}} \\
\vspace{2mm}
\includegraphics[width=1\textwidth]{image/certificate.png} \\
\begin{center}
    \large\textbf{Department of Computer Science \& Engineering}\\
    \large\textbf{(ARTIFICIAL INTELLIGENCE AND MACHINE LEARNING)}\\
    \small\textbf{Devarakaggalahalli, Harohalli,
Kanakapura Road,
Bengaluru South Dt. – 562 112\\}
    \small\textbf{Karnataka, India}
\end{center}
\vspace{5mm}
\Large{\underline{\textbf{CERTIFICATE}}} \\
  \end{center}
  \vspace{5mm}
\normalsize
This is to certify that the project entitled \textbf{“Crop Stress Grading using Hyperspectral Data and Deep Learning”} is carried out by \textbf{VANSH SONI}, 
bonafide student of Bachelor of Technology in Computer Science and Engineering at the School of Engineering, Dayananda Sagar University, Bangalore, in partial fulfillment for the award of a degree in Bachelor of Technology in Computer Science and Engineering, during the year 
\textbf{2025 \-- 2026}.\\

\vspace{20mm}
\noindent
\begin{minipage}[t]{0.31\textwidth}
\small
  \textbf{Prof. Udayabhaskara N}\\
  Assistant Professor\\
  Dept. of CSE (AI\&ML)\\
  School of Engineering\\
  Dayananda Sagar University\\
  \\
  Signature .........................\\
\end{minipage}%
\hfill
\begin{minipage}[t]{0.31\textwidth}
\small
  \textbf{Project Co-ordinator}\\
  Dept. of CSE (AI\&ML)\\
  School of Engineering\\
  Dayananda Sagar University\\
  \\
  Signature .........................\\
\end{minipage}%
\hfill
\begin{minipage}[t]{0.34\textwidth}
\small
  \textbf{Dr. Jayavrinda Vrindavanam}\\
  Professor \& Chairperson\\
  Dept. of CSE (AI\&ML)\\
  School of Engineering\\
  Dayananda Sagar University \\
  \\
  Signature .........................\\
\end{minipage}

\vspace{15mm}
\text{Name of the Examiners:}
\hfill
\vspace{5mm}
\text{Signature with date:} \\
\vspace{5mm}
\text{1...........................}
\hfill{.............................} \\
\vspace{5mm}
\text{2.............................} 
\hfill{.............................} \\


\newpage
  \pagestyle{plain}
\thisfancypage{%
  \setlength{\fboxsep}{20pt}\doublebox}{}
  \begin{center}
    \huge\textbf{DECLARATION}
\end{center}
\vspace{15mm}
I, \textbf{VANSH SONI}, am a student of B.Tech in Computer Science and Engineering (AI \& ML) at the School of Engineering, Dayananda Sagar University. I hereby declare that the Major Project titled \textbf{“Crop Stress Grading using Hyperspectral Data and Deep Learning”} has been carried out by me and submitted in partial fulfillment for the award of a degree in \textbf{Bachelor of Technology} in \textbf{Computer Science and Engineering} during the academic year \textbf{2025--2026}.
\vspace{15 mm}


\noindent
\begin{minipage}{0.8\textwidth}
    \textbf{Student:} \hfill\textbf{Signature}\\
    
    \textbf{Name:} VANSH SONI \\
\end{minipage}

\begin{flushleft}
    \textbf{Place:} Bangalore \\
    \textbf{Date:} 
\end{flushleft}


\newpage
  \pagestyle{fancy}
\thisfancypage{%
  \setlength{\fboxsep}{20pt}\doublebox}{}
\begin{center}
    \huge\textbf{ABSTRACT}
\end{center}
\vspace{15mm}
This study investigates the application of deep learning models for classifying crop stress stages using 1D hyperspectral data. Hyperspectral imaging provides a rich spectrum of data, capturing reflectance across numerous wavelengths, which is crucial for identifying early signs of stress in crops that are invisible to the naked eye or traditional multispectral cameras.

In this project, a robust data pipeline was developed to preprocess the hyperspectral data. This included missing value imputation, Savitzky-Golay filtering for spectrum smoothing, first-derivative spectroscopy, and the computation of targeted indices like the Moisture Stress Index (MLVI) and Hyperspectral Vegetation Stress Index (H-VSI). Data augmentation was performed using SMOTE to handle class imbalances across the six crop stress stages.

Four distinct deep neural network architectures were designed and evaluated: a 1D Convolutional Neural Network (CNN), a 1D ResNet, a Transformer model, and a Hybrid model combining Conv1D layers with a Transformer encoder. The models were trained using a cosine annealing learning rate scheduler and evaluated on a hold-out test set. The Hybrid model demonstrated the most promising results, achieving the highest accuracy of 90.10\%, effectively capturing both local spectral features and global dependencies. This research highlights the transformative potential of combining hyperspectral data with advanced AI techniques for early, automated, and accurate crop stress grading, paving the way for targeted agricultural interventions and enhanced crop yield.


\newpage
    \pagestyle{combined}
\thisfancypage{%
  \setlength{\fboxsep}{20pt}\doublebox}{}
  \pagenumbering{arabic}
\chapter{INTRODUCTION}
In the field of precision agriculture, ensuring crop health and optimizing yield are paramount challenges. Crops are frequently subjected to various biotic and abiotic stresses such as diseases, pest infestations, water scarcity, and nutrient deficiencies. The critical need to discover and grade these stresses early has been highlighted by the pressing need to begin treatments promptly, which can greatly affect agricultural outcomes and food security. The use of hyperspectral data as a diagnostic tool to reveal complex patterns of crop reflectance is the focus of this report, which aims to provide an automated deep learning solution.

\section{Context and Objective}
\vspace{2mm}
\begin{itemize}
    \item \textbf{Context:} Crop stress presents a significant obstacle in modern agriculture. The prompt diagnosis of stress stages is of utmost importance in order to initiate targeted interventions that have the potential to greatly influence crop yield and quality.
    \item \textbf{Objective:} This project examines the application of 1D Hyperspectral Data as a diagnostic modality for crop stress grading. The primary objective is to elucidate complex spectral reflectance patterns in order to facilitate early detection and enhance the accuracy of crop stage predictions using artificial intelligence.
    \item \textbf{Scope:} The focus is on addressing specific spectral anomalies and extracting relevant features, including computing first-derivative spectroscopy and specialized indices like MLVI and H-VSI. The objective is to evaluate multiple deep learning architectures (CNN, ResNet, Transformer, Hybrid) to clearly define distinct spectral patterns linked to various stress stages.
    \item \textbf{Significance:} Agricultural workers and automated systems can distinguish different stages of crop stress by recognizing and comprehending these spectral trends. The process of differentiation plays a crucial role in ensuring accuracy in resource planning and the prompt implementation of irrigation or fertilization techniques.
\end{itemize}


\newpage
  \pagestyle{fancy}
\thisfancypage{%
  \setlength{\fboxsep}{20pt}\doublebox}{}
\normalsize
\chapter{PROBLEM DEFINITION AND OBJECTIVE}
\section{Problem Definition}
\vspace{2mm}
Early detection of crop stress is a significant challenge in precision agriculture. Identifying stressed crops promptly is crucial for initiating therapy (such as targeted watering or pesticide application) that can significantly enhance yield outcomes. Currently, visual inspection is time-consuming and often detects stress only after irreversible damage has occurred. There exists a need for a precise and reliable diagnostic approach to detect crop stress at its earliest stages.

\begin{itemize}
    \item \textbf{Approach:} Hyperspectral data collection is a non-invasive method that is highly useful for diagnosing plant health by capturing patterns of spectral reflectance across hundreds of narrow contiguous wavelength bands.
    \item \textbf{Key Functions:} The utilization of hyperspectral data can facilitate the identification of abnormalities commonly linked with crop stress, such as changes in chlorophyll absorption, water content, and cellular structure. 
    \item \textbf{Distinguishing Patterns:} Through meticulous analysis of hyperspectral data, it is possible to identify specific patterns that are exclusive to different stages of crop growth and stress. 
\end{itemize}

\section{Novelty of Proposed Approach}
\vspace{2mm}
\begin{itemize}
    \item \textbf{High-Dimensional Data Processing:} Addresses the pressing issue of analyzing high-dimensional 1D hyperspectral signals, which contain hundreds of overlapping bands.
    \item \textbf{Advanced Preprocessing:} Employs Savitzky-Golay filtering for smoothing and computes the first derivative to highlight rapid changes in the spectral curve (e.g., the red edge).
    \item \textbf{Feature Engineering:} Integrates specific domain-knowledge indices, including Moisture Stress Index (MLVI) and Hyperspectral Vegetation Stress Index (H-VSI), utilizing specific Near-Infrared (NIR) and Short-Wave Infrared (SWIR) bands.
    \item \textbf{Deep Learning Models Employed:} Explores traditional 1D CNNs, Deep Residual Networks (ResNet1D), sequence-to-sequence Transformer architectures, and a novel Hybrid (Conv1D + Transformer) approach to capture both local spatial features and global contextual dependencies in the spectrum.
\end{itemize}

\newpage
  \pagestyle{fancy}
\thisfancypage{%
  \setlength{\fboxsep}{14pt}\doublebox}{}
\normalsize
\chapter{METHODOLOGY}
\section{Dataset Collection and Preprocessing}
\vspace{2mm}
To gather and prepare a diverse and comprehensive dataset for hyperspectral signal analysis, a robust data pipeline was established.

\begin{itemize}
    \item \textbf{Data Cleaning:} The initial dataset comprised spectral bands and metadata. Columns with more than 50\% missing values were identified and dropped. Rows containing residual missing values were subsequently removed to ensure data integrity.
    \item \textbf{Savitzky-Golay Filtering:} To eliminate high-frequency noise from the spectral signatures without distorting the signal tendency, a Savitzky-Golay filter was applied across the spectral axis.
    \item \textbf{Derivative Spectroscopy:} The first derivative of the smoothed spectral data was computed to emphasize the slopes of the reflectance curves, which are highly correlated with pigment concentration and structural changes.
    \item \textbf{Index Extraction:} Crucial bands were isolated—specifically NIR (approx. 854 nm), SWIR1 (approx. 1649 nm), and SWIR2 (approx. 2133 nm)—to compute the MLVI and H-VSI indices. These indices provide targeted insights into canopy water content and structural integrity.
    \item \textbf{Data Augmentation:} To address class imbalance among the six target stages, the Synthetic Minority Over-sampling Technique (SMOTE) was applied to the training set, ensuring the models do not become biased toward the majority class.
\end{itemize}

\section{Deep Learning Architectures}
\vspace{2mm}
Four distinct deep learning models were designed to classify the 1D hyperspectral sequences:
\begin{itemize}
    \item \textbf{Paper CNN:} A baseline 1D Convolutional Neural Network consisting of two convolutional layers with ReLU activations, followed by max pooling, dropout, and a dense classification head.
    \item \textbf{ResNet1D:} A deep residual network adapted for 1D signals. It features an initial convolutional stem followed by multiple stacked residual blocks (ResBlock1D) with increasing channel dimensions (up to 1024) and strided convolutions for downsampling, terminating in a global average pooling layer.
    \item \textbf{Transformer:} A sequence modeling architecture utilizing positional encoding and Multi-Head Self-Attention mechanisms to process the entire spectrum simultaneously, identifying long-range dependencies between distant spectral bands.
    \item \textbf{Hybrid Model:} This model first employs Conv1D layers (with residual skip connections) to extract local spectral textures and reduce the sequence length. The reduced, feature-rich sequence is then passed into a Transformer Encoder to map global inter-band relationships, ultimately feeding into a classification head.
\end{itemize}

\newpage
  \pagestyle{fancy}
\thisfancypage{%
  \setlength{\fboxsep}{20pt}\doublebox}{}
\normalsize
\chapter{REQUIREMENTS}
\section{Hardware Requirements}
\vspace{2mm}
\begin{itemize}
    \item \textbf{Processor:} Quad-core or higher processor (e.g., Intel Core i5, AMD Ryzen 5).
    \item \textbf{RAM:} 16GB or higher for efficient data processing and loading hyperspectral arrays.
    \item \textbf{Storage:} SSD with a minimum of 512GB for fast data access.
    \item \textbf{Graphics:} Dedicated GPU (NVIDIA GeForce GTX/RTX with CUDA support) is critical for accelerated training of deep learning models like ResNet and Transformers.
\end{itemize}

\section{Software Requirements}
\vspace{2mm}
\begin{itemize}
    \item \textbf{Python:} Version 3.8 or later.
    \item \textbf{Deep Learning Framework:} PyTorch for defining, training, and evaluating neural network architectures.
    \item \textbf{Data Processing Libraries:} Pandas and NumPy for tabular data manipulation and numerical operations. Scipy for signal processing (Savitzky-Golay filter).
    \item \textbf{Machine Learning Utilities:} Scikit-learn for data splitting, scaling, and computing evaluation metrics (Accuracy, F1-score, MCC). Imbalanced-learn for SMOTE data augmentation.
    \item \textbf{Visualization:} Tqdm for progress tracking during training loops.
\end{itemize}

\newpage
  \pagestyle{fancy}
\thisfancypage{%
  \setlength{\fboxsep}{20pt}\doublebox}{}
\normalsize
\chapter{EXPERIMENTATION AND RESULTS}
\section{Training Strategy}
\vspace{2mm}
The models were trained using a structured and rigorous methodology to ensure fair comparison and optimal performance.
\begin{itemize}
    \item \textbf{Loss Function:} Cross-Entropy Loss was utilized. To further handle class imbalance, class weights were computed based on inverse class frequencies and passed to the loss function. Label smoothing (0.1) was applied to prevent overconfidence and improve generalization.
    \item \textbf{Optimizer \& Scheduler:} The AdamW optimizer was used with an initial learning rate of 1e-4 and a weight decay of 1e-4. A Cosine Annealing Warm Restarts scheduler (T0=10, T\_mult=2) was employed to dynamically adjust the learning rate and help the models escape local minima.
    \item \textbf{Training Loop:} Models were trained for up to 300 epochs with a batch size of 64. Early stopping was implemented with a patience of 15 epochs, monitoring the Matthews Correlation Coefficient (MCC) on the validation set to halt training when performance plateaued.
    \item \textbf{Gradient Clipping:} Gradient clipping (max\_norm=1.0) was applied to stabilize training, particularly for the Transformer and Hybrid models, preventing explosive gradients.
\end{itemize}

\section{Performance Analysis}
\vspace{2mm}
The models were evaluated on a completely unseen test set. The primary metrics used for evaluation were Accuracy, F1-Score, and the Matthews Correlation Coefficient (MCC).

\begin{table}[H]
    \centering
    \begin{tabular}{|c|c|c|}
    \hline
     \textbf{DEEP LEARNING MODEL} & \textbf{ACCURACY} & \textbf{MCC} \\
   \hline
   Paper CNN & 83.27 \% & 0.7931 \\
   \hline
   ResNet1D & 89.23 \% & 0.8662 \\
    \hline
   Transformer & 89.04 \% & 0.8637 \\
    \hline
   \textbf{Hybrid (CNN + Transformer)} & \textbf{90.10 \%} & \textbf{0.8767} \\
    \hline
    \end{tabular}
    \caption{Performance metrics of different deep learning models on the test set}
    \label{tab:results}
\end{table}

The Hybrid model demonstrated the superior capability, achieving the highest accuracy of 90.10\%. This indicates that combining convolutional layers for local spectral feature extraction with a Transformer encoder for global dependency modeling provides the optimal inductive bias for hyperspectral crop data. ResNet1D also performed exceptionally well, closely followed by the standalone Transformer, while the simpler Paper CNN lagged behind.

\newpage
  \pagestyle{fancy}
\thisfancypage{%
  \setlength{\fboxsep}{20pt}\doublebox}{}
\normalsize
\chapter{CONCLUSION AND FUTURE SCOPE}

\section{Conclusion}
\vspace{2mm}
This project successfully demonstrated the application of advanced deep learning techniques on 1D hyperspectral data for crop stress grading. By establishing a robust preprocessing pipeline—incorporating Savitzky-Golay smoothing, derivative spectroscopy, and vegetation indices (MLVI, H-VSI)—the spectral data was transformed into a highly informative feature space. 

The empirical evaluation of four distinct neural network architectures (CNN, ResNet1D, Transformer, and Hybrid) provided valuable insights. The Hybrid model, leveraging the strengths of both convolutions and self-attention mechanisms, emerged as the most effective approach, yielding a peak accuracy of 90.10\% across the six stress stages. This represents a significant step forward in automated agricultural diagnostics, proving that complex spectral reflectance patterns can be reliably decoded by artificial intelligence to detect crop stress.

\section{Future Scope}
\vspace{4mm}
\begin{itemize}
    \item \textbf{Deployment and Edge Computing:} The next logical step is deploying the trained Hybrid or ResNet models onto edge devices, such as Unmanned Aerial Vehicles (UAVs) equipped with hyperspectral sensors, enabling real-time, on-the-fly crop stress mapping.
    \item \textbf{Multimodal Integration:} Integrating the hyperspectral data pipeline with other modalities, such as thermal imaging, LiDAR structural data, or weather patterns, could further enhance diagnostic accuracy and robustness.
    \item \textbf{Explainable AI (XAI):} Applying techniques like Integrated Gradients or Attention Rollout to the Transformer components to understand exactly which spectral bands are driving the model's predictions. This interpretability can help agronomists discover new spectral biomarkers for specific crop diseases.
    \item \textbf{Dataset Expansion:} Expanding the dataset to include varying geographical regions, different crop types, and diverse lighting conditions will be crucial for ensuring the model's generalization capabilities in real-world, large-scale agricultural environments.
\end{itemize}

\end{document}
